% Requirements in preamble:
% \usepackage{graphicx}
% \usepackage{booktabs}
% \usepackage{svg} % for \includesvg{...}
% Compile with --shell-escape when using \includesvg.

\section{Resultats}

\subsection{Vue d'ensemble}
L'evaluation montre une recuperation de cle AES tres solide sur les jeux de traces alignes, avec une recuperation complete (16/16 octets au rang 1) sur plusieurs scenarios a 20k traces. Les experiments en mode blind strict et en configuration highside deux-passes restent en revanche nettement plus difficiles, ce qui fournit une limite experimentale importante.

\subsection{Tableau des resultats clefs}
\begin{table}[h]
\centering
\begin{tabular}{lrrrr}
\toprule
Experience & Traces & Rank-1 (sur 16) & Top-5 (sur 16) & Rang moyen \\
\midrule
\texttt{unknown\_B\_5k\_aligned\_fullscan} & 5000 & 14 & 16 & 1.1875 \\
\texttt{unknown\_B\_20k\_aligned\_fullscan} & 20000 & 16 & 16 & 1.0000 \\
\texttt{pres\_unknown\_blind\_20k} & 20000 & 15 & 16 & 1.0625 \\
\texttt{pres\_unknown\_B\_20k} & 20000 & 16 & 16 & 1.0000 \\
\texttt{pres\_unknown\_new\_20k} & 20000 & 16 & 16 & 1.0000 \\
\texttt{highside\_two\_pass} (pass1) & 1000 & 1 & 3 & 85.8750 \\
\texttt{highside\_two\_pass} (pass2) & 1000 & 1 & 1 & 153.5000 \\
\texttt{strict\_blind\_B\_cfg2pp\_ensemble} (mean\_score) & 15000 & 0 & 0 & 158.4375 \\
\bottomrule
\end{tabular}
\caption{Resultats quantitatifs principaux.}
\label{tab:resultats-clefs}
\end{table}

\subsection{Interpretation}
\begin{enumerate}
\item \textbf{Resultat principal}: la cle est recuperee completement (16/16 rang 1) sur \texttt{unknown\_B\_20k\_aligned\_fullscan} et sur les deux scenarios de presentation (\texttt{pres\_unknown\_B\_20k}, \texttt{pres\_unknown\_new\_20k}).
\item \textbf{Effet du nombre de traces}: on observe un gain net entre 5k et 20k traces (14/16 $\rightarrow$ 16/16 en rank-1) sur la meme famille de donnees \texttt{unknown\_B}.
\item \textbf{Robustesse en blind}: \texttt{pres\_unknown\_blind\_20k} reste tres performant (15/16 en rank-1, 16/16 en top-5), ce qui indique une attaque exploitable meme sans connaissance parfaite.
\item \textbf{Limites}: \texttt{highside\_two\_pass} et \texttt{strict\_blind} confirment qu'en conditions plus dures, la methode actuelle ne converge pas correctement vers la cle complete.
\end{enumerate}

\subsection{Figures (SVG)}

\begin{figure}[h]
\centering
\includesvg[width=0.95\linewidth]{figures/convergence_aes_2204_20k}
\caption{Convergence CPA - campagne \texttt{aes\_2204}.}
\label{fig:conv-aes2204}
\end{figure}

\noindent
\textit{Lecture}: l'augmentation du nombre de traces fait monter rapidement le nombre d'octets recuperes en rank-1 et atteint 15/16 a 10k traces, avec 16/16 en top-5.

\begin{figure}[h]
\centering
\includesvg[width=0.95\linewidth]{figures/convergence_aes_step_16b_20k}
\caption{Convergence CPA - campagne \texttt{aes\_step\_16b}.}
\label{fig:conv-aesstep16b}
\end{figure}

\noindent
\textit{Lecture}: la convergence existe mais est plus lente; le passage au voisinage de 15k-20k traces est necessaire pour stabiliser les octets au rang 1.

\begin{figure}[h]
\centering
\includesvg[width=0.95\linewidth]{figures/comparaison_resultats_principaux}
\caption{Comparaison des resultats principaux.}
\label{fig:comparaison-principaux}
\end{figure}

\noindent
\textit{Lecture}: ce graphe resume l'ecart entre scenario 5k et 20k, et montre que le mode blind 20k reste tres proche du scenario non blind.

\begin{figure}[h]
\centering
\includesvg[width=0.95\linewidth]{figures/resultats_limites}
\caption{Resultats limites (negatifs).}
\label{fig:resultats-limites}
\end{figure}

\noindent
\textit{Lecture}: le rang moyen eleve en \texttt{highside\_two\_pass} (surtout pass2) et en \texttt{strict\_blind} met en evidence les conditions ou l'attaque n'est pas encore fiable.

\subsection{Message a garder pour la conclusion}
La methode est \textbf{efficace et reproductible} sur des jeux alignes avec un budget de traces suffisant (20k), tout en gardant une bonne performance en blind. En revanche, les pipelines highside/deux-passes et blind strict constituent encore des \textbf{verrous de robustesse} a traiter pour generaliser l'attaque.
